\section{Uniform computation and scalar-image constraints}\label{sec:computation}

We represent a finite computation by a circuit whose exponential address
space has a polynomial-size description, and then express the absence of
accepting tests by scalar-image consistency. The computation proposition
is stated for an arbitrary fixed deterministic machine: choosing a
complement decider is a separate step at the end of this section.
The product-address and buffer methods also appear in the fixed-dimension
companion \citep[Sections~3--4]{CompanionFixedDimension}; here every bound
is uniform in the growing number of binary coordinates.

\paragraph{Circuit conventions.}
A scalar gate takes the OR of its incoming connections and may additionally
be declared a \emph{seed}. It is true precisely when it is a seed or at
least one incoming source is true. An AND gate has two input ports; each
port takes the OR of its incoming connections, and the gate is true when
both ports are true. Empty ORs are false. The circuit is finite and
acyclic, so these rules determine all truth values. Designated
\emph{tests} are scalar gates with no outgoing connections.

\begin{definition}[Product description]\label{def:computation-product}
For an address length $r$, every gate type has one gate at each address in
$\{0,1\}^r$. A connection schema specifies source and destination types,
a destination port, nonempty coordinate domains
$A_i\subseteq\{0,1\}$, and injections
$\sigma_i:A_i\longrightarrow\{0,1\}$, for $1\le i\le r$.
It describes the connections
\[
 (\text{source type},\mathbf a)
 \longrightarrow
 (\text{destination type},\boldsymbol\sigma(\mathbf a)),
 \qquad
 \mathbf a\in\prod_{i=1}^r A_i,
\]
at the specified port, where
$\boldsymbol\sigma(\mathbf a)=(\sigma_i(a_i))_{i=1}^r$.
A seed or test request specifies a scalar type and a product domain of
addresses. The description lists these schemas and requests by their
one-coordinate tables, rather than listing all addressed gates.
\end{definition}

Let \(M\) be a fixed deterministic one-tape machine on a doubly infinite
tape, with a finite alphabet and head displacements in \(\{-1,0,1\}\).
Its head starts at position zero, the input \(w\) of length \(m\) occupies
positions \(0,\ldots,m-1\), and all other cells are blank. Halting
configurations are frozen. Let \(p\) be a fixed integer polynomial with
\(p(m)\ge m+2\) such that \(M\) halts on every input of length \(m\)
within \(T=2^{p(m)}\) steps. Put
\begin{equation}\label{eq:computation-parameters}
 d=p(m)+2,\qquad H_*=2^d=4T,\qquad r=2d,\qquad K=r+1.
\end{equation}
Thus \(r\ge8\). An address is a pair \((t,j)\) of \(d\)-bit numbers
with \(0\le t,j<H_*\). Time is not cyclic; position is interpreted
modulo \(H_*\).

\Needspace{18\baselineskip}
\begin{proposition}[Uniform product-circuit compilation]\label{prop:computation}
Fix a deterministic one-tape machine \(M\) with the initialization and
frozen-halting convention just specified, and an integer polynomial \(p\)
satisfying \(p(m)\ge m+2\) and bounding its running time by \(2^{p(m)}\).
For every word \(w\) of length \(m\), one can generate a product description
in the sense of \Cref{def:computation-product} of a finite acyclic circuit
at address length \(r=2(p(m)+2)\), with scalar terminal tests, such that
at least one test is true if and only if \(M\) accepts \(w\).

There are \(O_M(1)\) gate types and \(O_M((r+1)^2+m)\) connection
schemas and seed or test requests, each with \(r\) constant-size
coordinate entries. All domains have nonempty binary factors and all
coordinate maps are injections. The total number of coordinate entries
is \(O_M(r((r+1)^2+m))\). The description is generated deterministically
in time polynomial in \(r+m\), with constants depending only on \(M,p\),
without enumerating times or full addresses.
\end{proposition}

\begin{proof}
Let $\Gamma$ be the fixed tape alphabet and let $\Sigma$ consist of its
plain symbols together with symbols carrying a head state of $M$. There
is a fixed map
\[
 F:\Sigma^3\longrightarrow\Sigma
\]
that updates the center cell of every configuration with exactly one head.
If the head is at the center, its state and scanned symbol determine the
symbol written there and whether the head remains. If it is at a neighbor,
that neighbor's marked symbol determines whether the head moves into the
center and, if so, its new state; the center's tape symbol is already in
the triple. If the triple has no head, the center stays unchanged. These
rules also cover halted states, for which nothing changes. Complete $F$
arbitrarily on triples impossible in a one-head configuration.

The same rule implements one machine step on a ring of length $H_*$.
Starting from one head, it writes at that head and transfers the head by
the prescribed displacement, so the ring still has exactly one head.
In particular, a triple crossing the ring's seam is a prescribed valid
local case, not an arbitrary entry of the table.

Initialize the ring with $w$ at $0,\ldots,m-1$, the head at zero, and
blanks elsewhere; for $m=0$ the head carries a blank. This simulates the
line computation through halting. All input positions and all positions
visited within $T$ steps lie in $[-T,T]$, since $m\le T$. Reduction modulo
$4T$ is injective on this interval, and the initial contents at these
positions agree with the corresponding ring cells. Inductively, each
scanned ring cell has received exactly the same previous writes as its
unique potentially visited line counterpart. The next transition
therefore writes the same symbol and moves to the same corresponding
position, including a last move to either endpoint of the interval.
There can be no change elsewhere: a valid machine update changes only
the old and new head cells. Equivalently, all neighborhoods that could
react to the head lie in the halo $[-T-2,T+2]$, whose $2T+5$ positions are
distinct modulo $4T$ because $T\ge4$; centers with no head in their triple
stay unchanged. Upon halting the whole ring configuration freezes.
Its configuration at time $H_*-1$ consequently records the outcome of $M$.

\paragraph{The gates.}
For each $f\in\Sigma$, introduce a scalar type $S_f$. Seed exactly
$S_f(0,j)$ when $f$ is the actual initial letter at position $j$, including
the head information. For every $(a,b,c)\in\Sigma^3$, introduce two AND
types $U_{abc}$ and $V_{abc}$. For $t>0$ use the connections
\begin{align}
 S_a(t-1,j-1)&\longrightarrow U_{abc}(t,j)\text{, port }1,
 \label{eq:computation-left}\\
 S_b(t-1,j)&\longrightarrow U_{abc}(t,j)\text{, port }2,
 \label{eq:computation-center}\\
 S_c(t-1,j+1)&\longrightarrow V_{abc}(t,j)\text{, port }2.
 \label{eq:computation-right}
\end{align}
At every time, including zero, also use
\begin{equation}\label{eq:computation-same-time}
 U_{abc}(t,j)\longrightarrow V_{abc}(t,j)\text{, port }1,
 \qquad
 V_{abc}(t,j)\longrightarrow S_{F(a,b,c)}(t,j).
\end{equation}
All position indices in these formulas are cyclic.

Assign ranks $3t$, $3t+1$, and $3t+2$ to the $U$, $V$, and $S$ gates at
time $t$, respectively. Every connection strictly increases rank, proving
acyclicity. At time zero, both ports of every $U$ are false, and every $V$
is false; precisely the intended scalar seeds are true. Suppose the true
scalars at time $t-1$ specify the actual letter at every cell. At $(t,j)$,
$U_{abc}$ is true exactly when $a,b$ are the actual left and center
letters, and $V_{abc}$ is true exactly when $c$ is also the actual right
letter. Exactly one triple makes $V$ true there, so exactly the scalar
for the next letter is true. Induction proves that the scalar gates
record the ring computation at every time.

For each letter $f$ carrying an accepting state, test all
$S_f(H_*-1,j)$. Some test is true exactly when $M$ accepts $w$.
These tests are terminal: a scalar's outgoing connections advance time,
and none starts at time $H_*-1$. The same-time connections
\eqref{eq:computation-same-time} do not leave scalar gates.

\paragraph{Carry and borrow tables.}
Number bits within either $d$-bit group from $0$ to $d-1$, least
significant first. For $0\le h<d$, define the product domains
\begin{align*}
 I_h^+&=\{x\in\{0,1\}^d:x_i=1\ (i<h),\ x_h=0\},\\
 I_h^-&=\{x\in\{0,1\}^d:x_i=0\ (i<h),\ x_h=1\}.
\end{align*}
Bits above $h$ are unrestricted. On either domain use the componentwise
map
\begin{equation}\label{eq:computation-bit-flip}
 \theta_{h,i}(x_i)=
 \begin{cases}
  1-x_i,&i\le h,\\
  x_i,&i>h.
 \end{cases}
\end{equation}
Its restriction to each coordinate domain is injective. The $d$ domains
$I_h^+$ partition all numbers below $H_*-1$, and their maps implement
addition of one. They therefore give time increment without overflow.
For cyclic position increment, add the all-ones singleton and flip all
bits on that additional domain. For cyclic decrement, the $I_h^-$
domains partition the nonzero numbers; their maps subtract one. Add the
all-zero singleton, again flipping all bits, for its wraparound case.
Thus each cyclic position shift has $d+1$ product cases.

In source coordinates, the temporal connections
\eqref{eq:computation-left}--\eqref{eq:computation-right} are, respectively,
\begin{equation}\label{eq:computation-source-maps}
 (t,j)\mapsto(t+1,j+1),\qquad
 (t,j)\mapsto(t+1,j),\qquad
 (t,j)\mapsto(t+1,j-1).
\end{equation}
Combining an $I_h^+$ time case with a position case remains a product
domain because the two digit groups are disjoint. The two shifted kinds
need $d(d+1)$ schemas each; the unshifted kind needs $d$. There is no
all-ones time case. Each connection in
\eqref{eq:computation-same-time} uses one identity schema.

\paragraph{Initial and final requests.}
Put $m'=\max\{1,m\}$, so $m'<H_*$. At time zero, give one singleton
position request for each $0\le j<m'$, using its actual initial letter.
When $m=0$ this specifies the head-bearing blank at position zero.
All remaining initial positions carry the plain blank. Write the
$d$ bits of $m'$ as $\mu_0,\ldots,\mu_{d-1}$. A disjoint product
decomposition of $j\ge m'$ consists of the singleton $\{m'\}$ and,
for every $h$ with $\mu_h=0$, the domain
\begin{equation}\label{eq:computation-blank-domain}
 j_i=\mu_i\quad(i>h),\qquad j_h=1,\qquad
 j_i\in\{0,1\}\quad(i<h).
\end{equation}
Indeed, for $j>m'$ the largest bit at which the two numbers differ has
exactly this form. Fixing every time bit to zero makes each such blank
request a product in all $r$ coordinates. There are at most $d+1$ blank
domains and hence at most $m'+d+1$ seed requests in total. Each accepting
letter needs one test request, fixing every time bit to one and leaving
the position bits unrestricted.

Finally, write $q=|\Sigma|$. There are $q+2q^3$ gate types. The connection
schemas number
\begin{equation}\label{eq:computation-schema-count}
 q^3\bigl(2d(d+1)+d+2\bigr)=q^3(2d^2+3d+2).
\end{equation}
The seed and test requests add $O_M(d+m+1)$ domains. Their total tables
therefore have $O_M(r((r+1)^2+m))$ constant-size coordinate entries.
To generate them, loop over the fixed alphabet, the carry or borrow
indices, the $m'$ initial positions, and the zero-bit positions of $m'$;
fill each coordinate using the displayed identity, flip, or fixed-bit
rule. The initial letters come directly from $w$. All loop lengths are
polynomial in $r+m$, and the counters and small position numbers have
$O_M(\log(r+m+2))$ bits. No loop runs through $H_*$ times or $2^r$
addresses. The bit-cost accounting, including evaluation of the fixed
polynomial $p$, is given in Proposition~\ref{prop:uniform-general}.
\end{proof}

\subsection{Nonempty sets and individual scalar images}\label{sec:consistency}

All vector spaces in the remaining construction are over $\F$. The
intermediate problem compares images of nonempty sets under individual
scalar forms. It does not ask for a simultaneous assignment of one vector
to every variable.

\begin{definition}[Product consistency instance]
\label{def:scalar-images}
Fix an address length $r$. A finite nonempty set $\mathcal P$ of
\emph{real types} has dimensions $d_P\in\{1,2\}$. For every $P\in\mathcal P$
and every $\mathbf a\in\{0,1\}^r$, there is a variable with value space
$\F^{d_P}$. A schema $e\in\mathcal E$ consists of ordered endpoint types
$P,R$, a product domain
\[
 A_e=\prod_{i=1}^r A_{e,i},\qquad
 \varnothing\ne A_{e,i}\subseteq\{0,1\},
\]
coordinate injections $\sigma_{e,i}:A_{e,i}\to\{0,1\}$, and scalar
affine forms $\phi_{e,0}$ on $\F^{d_P}$ and $\phi_{e,1}$ on $\F^{d_R}$.
Write $\boldsymbol\sigma_e$ for their componentwise address map.
Endpoint types may coincide, parallel schemas are allowed, and the
linear part of an affine form may be zero.

The instance is \emph{consistent} if there are nonempty sets
$Q_{P,\mathbf a}\subseteq\F^{d_P}$ for all variables such that
\begin{equation}\label{eq:scalar-images}
 \phi_{e,0}(Q_{P,\mathbf a})
 =\phi_{e,1}(Q_{R,\boldsymbol\sigma_e(\mathbf a)})
 \qquad(e\in\mathcal E,\ \mathbf a\in A_e).
\end{equation}
Here $\phi(Q)=\{\phi(u):u\in Q\}$. Each equality compares one scalar
image at each end; several comparisons are never combined into a joint
image condition.
\end{definition}

This is a per-comparison local-support condition in the classical
consistency tradition \citep[Sections~4 and~6]{Mackworth1977}:
every value at one endpoint has some support at the other, in both
directions. Supports can differ between comparisons and between the
two occurrences of a self-comparison. This viewpoint does not assert
one globally satisfying vector at every variable.
Related existential pebble-game arguments use extendible families of
partial maps \citep[Definition~4]{Berkholz2013}. Our later bijective-game
argument must additionally select a whole bijection \emph{before}
Spoiler chooses the new vertex; exact projection of affine offsets
will supply that stronger order of choices.

The buffer below implements directed Boolean propagation using symmetric
image equalities. One-way propagation already plays a central role in
Grohe's monotone-circuit reductions for finite-variable equivalence
\citep[author manuscript, Lemma~14 and Section~5.4]{Grohe1999}
and in the later identification
construction of Lichter, Ra{\ss}mann and Schweitzer
\citep[Section~5]{LichterRassmannSchweitzer2025}.
Our private two-coordinate scalar-image buffer implements that directional
effect by a different construction, proved directly below; neither a
graph-switch theorem nor a growing-dimension bound is imported.
Its two coordinates may use different witnesses from
the source set; this is exactly why it transmits a forced zero without
turning the intermediate problem into a global vector assignment.

\Needspace{5\baselineskip}
\begin{proposition}[Circuit consistency]\label{prop:consistency}
For every finite acyclic circuit with the conventions of
Section~\ref{sec:computation}, one can construct nonempty scalar-image
constraints with vector dimensions at most two which are consistent if
and only if every test is false.

If the circuit has a product description with $g$ gate types, $c$
connection schemas, and $b$ seed and test requests in total, the result is
an instance of Definition~\ref{def:scalar-images} with $g+c+1$
real types and $3c+b$ comparison schemas, at the same address length.
Its one-coordinate domains and maps are copied from that description or
are identities. In particular, the circuit in
Proposition~\ref{prop:computation} gives
$O_M((r+1)^2+m)$ real types and schemas with dimensions in $\{1,2\}$.
\end{proposition}

\begin{proof}
We first describe the constraints for the expanded finite circuit, then
give their product representation. A scalar gate has one scalar variable;
its input and output directions are both the identity. An AND gate has
a vector variable $(x_1,x_2)\in\F^2$. Its two input directions are the
coordinate projections, and its output direction is $x_1+x_2$.

Introduce one dummy scalar variable. For each seed, compare its scalar
direction with the constant-zero form on the dummy. For each test, compare
its scalar direction with the constant-one form on the dummy. The dummy
may have any nonempty set of values: its two constant forms have images
$\{0\}$ and $\{1\}$, respectively, without any restriction on that set.

For each individual connection, with source output direction $o$ and
destination input direction $s$, introduce a private buffer
$(z,z')\in\F^2$ and impose the three separate image comparisons denoted by
\begin{equation}\label{eq:consistency-buffer}
 o=z,\qquad o=z',\qquad z+z'=s.
\end{equation}
The first two equalities do not require a common source value to be
simultaneously equal to both buffer coordinates.

\paragraph{A true test prevents consistency.}
Consider any proposed family of nonempty sets satisfying all comparisons.
If a source output image is $\{0\}$, the first two comparisons in
\eqref{eq:consistency-buffer} force every buffer vector to have both
coordinates zero. Nonemptiness makes its set exactly $\{(0,0)\}$, so
the third comparison forces the destination direction to have image
$\{0\}$.

Induct in a topological order of the circuit. A true scalar gate is either
seeded, in which case its constant comparison gives image $\{0\}$, or
has a true predecessor, whose zero output image propagates to it. At a
true AND gate each port has a true incoming source; both coordinate images
are therefore $\{0\}$, and the sum output also has image $\{0\}$.
Every true gate consequently has zero output image under every consistent
family. A true test would also require image $\{1\}$, a contradiction.

\paragraph{Witness sets when all tests are false.}
Assign to each test scalar the set $\{1\}$. At every other scalar gate
use $\{0\}$ if it is true and $\F$ if it is false. At an AND gate use
\begin{equation}\label{eq:consistency-and-witness}
 \{(x_1,x_2)\in\F^2:
       x_j=0\text{ for each true input port }j\}.
\end{equation}
All these sets are nonempty. A true AND has zero output image; a false
AND has at least one unconstrained coordinate, so its sum output image
is $\F$. Thus every true source has output image $\{0\}$ and every
false source has output image $\F$: the exceptional false tests are
terminal and never serve as sources.

For a connection with true source, the destination port is true, and its
chosen direction image is $\{0\}$. Use the buffer set $\{(0,0)\}$.
For a connection with false source, let $D\subseteq\F$ be the already
chosen nonempty image of its destination direction and use
\begin{equation}\label{eq:consistency-fiber}
 B_D=\{(z,z')\in\F^2:z+z'\in D\}.
\end{equation}
Its sum image is exactly $D$, and each coordinate image is $\F$:
given $d\in D$, the pairs $(u,u+d)$ as $u$ runs over $\F$ already
give both full coordinate images. Explicitly,
\[
\begin{array}{c|c|c}
 D & B_D & \text{each coordinate image}\\ \hline
 \{0\} & \{(0,0),(1,1)\} & \F\\
 \{1\} & \{(0,1),(1,0)\} & \F\\
 \F & \F^2 & \F
\end{array}
\]
Hence all three comparisons hold. In particular, a false source can feed
a zero-forced port through the diagonal buffer, or a false test through
the off-diagonal buffer, while retaining its full output image. Different
incoming or outgoing connections have private buffers and require only
these individual image equalities, so they impose no further coupling.
The seed and test comparisons also hold; take the dummy set to be $\F$.
This constructs a consistent family.

\paragraph{Preserving the product description.}
Give each gate type its corresponding one- or two-dimensional real type,
with variables at all addresses. Give each connection schema its own
two-dimensional buffer type, indexed in the source address system. On its
source product domain, the first two comparisons in
\eqref{eq:consistency-buffer} use identity address transmission from the
source gate to its buffer. The third uses the original componentwise map
from that buffer to the destination gate. Thus an allowed buffer address
belongs to exactly the intended individual connection of that schema.
Overlapping domains of different schemas use different buffer types.
Buffer variables outside their schema's domain are unused and may receive
the full set $\F^2$.

Use one scalar dummy type with variables at all addresses. Seed and test
requests become identity-transmission comparison schemas on the same
product domains, with the appropriate constant form at the dummy end.
This gives exactly one additional type per connection schema, one dummy
type, three comparisons per connection schema, and one comparison per
request. Every coordinate factor remains nonempty and every transmitted
map is an injection. The displayed count follows, and the tables can be
produced directly without expanding the addressed circuit.
\end{proof}

\begin{corollary}[The polarity of the reduction]\label{prop:computation-images}
For every binary language \(L_0\in\EXPTIME\), there are a fixed machine
\(M\) and fixed integer polynomial \(p\), with \(p(m)\ge m+2\), such that
a word \(w\) of length \(m\) can be transformed into a product
scalar-image instance at \(r=2(p(m)+2)\ge8\) that is consistent exactly
when \(w\in L_0\). It has \(O_M((r+1)^2+m)\) real types and
comparison schemas, with dimensions in \(\{1,2\}\), and is generated
in deterministic polynomial time in \(m\).
\end{corollary}
\begin{proof}
Choose an exponential-time decider for the complement of \(L_0\),
by swapping the two halting outputs of a decider for \(L_0\).
A fixed multitape machine can be simulated on one tape with polynomial
overhead: store its tapes and head markers on fixed tracks, and scan the
used interval to collect scanned symbols and then perform the updates.
After \(t\) simulated steps the interval has length \(O_M(m+t+1)\);
a fixed number of scans per step gives polynomial total overhead.
Enlarging a fixed integer polynomial \(p\) absorbs that overhead and
ensures \(p(m)\ge m+2\). Apply \Cref{prop:computation} to this normalized
complement decider, then \Cref{prop:consistency}. Consistency is equivalent
to all tests being false, hence to \(M\) not accepting \(w\), which is
equivalent to \(w\in L_0\). The exact buffer counts give the claimed
description size. Since \(r\) is polynomial in \(m\) for fixed \(p\),
the generation time is polynomial in \(m\) as well.
\end{proof}
